Sparse allocation does not eliminate work on particles throughout a filled sand bed. Both full-terrain methods increase particle count from 972,000 at 3\,m to 10,800,000 at 10\,m. The following table compares the largest size completed by both methods for each scene. Times are mean $\pm$ standard deviation in milliseconds; memory is the measured mean in GiB.
\begin{center}\small\begin{tabular}{lrrrrr}\toprule
Scene & Side [m] & Fixed [ms] & Sparse [ms] & Fixed [GiB] & Sparse [GiB] \\\midrule
Box & 8 & $80.96\pm6.23$ & $186.86\pm2.68$ & 12.87 & 22.98 \\
Rover & 7 & $67.18\pm6.65$ & $148.46\pm3.17$ & 14.41 & 22.63 \\
G1 & 7 & $66.05\pm1.66$ & $142.98\pm2.39$ & 14.33 & 22.41 \\
\bottomrule\end{tabular}\end{center}
Sparse has a lower mean step time in 0 of 16 matched size/scene comparisons. These results describe the tested reserved-capacity configuration; they do not establish an inherent advantage for either grid representation. No component-level profiling was performed to separate grid rebuilding, material solving, and capacity-related costs.
5 attempts failed during GPU memory allocation with the tested settings.
4 initialization workers were terminated by SIGKILL before measurements; the cause could not be confirmed because kernel logs were not accessible.
Failed measurements are excluded from means and shown as missing plot points. The affected cases and failure messages are listed in the accompanying summary data.
Box sparse: failed runs at side lengths 10, 9\,m.
Rover fixed: failed runs at side lengths 10\,m.
Rover sparse: failed runs at side lengths 10, 8, 9\,m.
G1 fixed: failed runs at side lengths 10\,m.
G1 sparse: failed runs at side lengths 10, 8, 9\,m.
